\documentclass[twocolumn,11pt]{article}
\usepackage[margin=0.75in]{geometry}
\usepackage{booktabs}
\usepackage{graphicx}

\graphicspath{{figures/}}

\title{Two-Column Layout Stress Case}
\author{Test Corpus}
\date{}

\begin{document}
\maketitle

\section{Overview}

This source uses the two-column article option. It includes a wide figure float
that spans both columns and a normal table that remains within one column. The
text around the floats is long enough to force ordinary column balancing.

The first paragraph describes a hypothetical evaluation setting in which sites
report weekly process indicators. The paragraph provides running text before the
wide figure so that the placement algorithm has realistic surrounding content.

\begin{figure*}[t]
\centering
\includegraphics[width=0.82\textwidth]{wide_figure.png}
\caption{A generated wide figure spanning both columns}
\label{fig:wide}
\end{figure*}

Additional text follows the wide display. The purpose is to create enough
material for column flow before and after the float. Converters often linearize
two-column layouts, so this test records how the result handles a document whose
source was explicitly typeset in columns.

\begin{table}[htbp]
\centering
\caption{Column-local summary}
\label{tab:twocol}
\begin{tabular}{lrr}
\toprule
Group & Mean & N \\
\midrule
North & 2.41 & 38 \\
South & 2.87 & 42 \\
East & 2.66 & 35 \\
West & 2.52 & 39 \\
\bottomrule
\end{tabular}
\end{table}

\section{Discussion}

The remaining paragraphs give the two-column layout enough content to exercise
pagination. Site-level summaries are intentionally simple, allowing the layout
features rather than the substantive claims to drive the test behavior. The table
in this section should remain readable even though the surrounding document uses
columns.

The wide figure in Figure \ref{fig:wide} is a separate case from the local table
in Table \ref{tab:twocol}. Both floats should have captions, labels, and nearby
paragraph text in the compiled PDF.

\end{document}
